\documentclass{article}
\usepackage{graphicx} % Required for inserting images
\usepackage{hyperref}
\usepackage{xcolor}

% Problems (numbered within each subsection)
\newcounter{problem}[subsection]
\renewcommand{\theproblem}{\thesubsection.\arabic{problem}}
\newcommand{\problem}[1]{%
  \refstepcounter{problem}%
  \par\noindent\textcolor{red}{\textbf{Problem \theproblem.}}\;#1\par
}

\newcommand{\vs}{
    \vspace{1em}
}

\newcommand{\vsn}[1]{
    \vspace{#1}
}
\title{High Performance Computing Notes}
\author{Kaustubh Singh}

% Make plans to port all this to a local platform.
% Github would serve well. 
% Check how to link this project to GitHub, 
% That way you can work locally, on Overleaf and sync changes via GitHub.

\begin{document}

\maketitle

\begin{center}
    As I write this, I have a solid background in the fundamentals of physics at the graduate level, but not so much in computing. I have taken some courses in scientific computing, which I found interesting and wanted to explore further, hence this endeavour. I am well-versed in Python, having created a library for hybrid tensor networks, but am new to principles of computing and compiled languages. 

    \vs 
    
    These notes will serve as reference for all that I learn as I teach myself high performance computing. I hope to sniff several rabbit holes, if not dive into them, as I begin learning C, a language closer to the hardware than Python. I will not shy away from being verbose. The language is meant to serve as a stepping stone as I muck around the muddy ponds of low-level, performant and parallel classical computing. To that end, I might pick up some different languages, but that will be way down the line. 
\end{center}

\section{Data types}
\subsection{Floating point arithmetic}
A succinct \href{https://www.phys.uconn.edu/~rozman/Courses/P2200_15F/downloads/floating-point-guide-2015-10-1}{article on floating point arithmetic}.

\vs 

\problem{Write a C program that outputs the bit sequence of a given data type. For instance, for float, you should show the signed bit, the exponent bits, the significand/ mantissa bits. Bonus points for nice formatting.}

\vs




\end{document}
